% concatenated from main.tex, Sec1_introduction.tex, Sec2_notation.tex, Sec3_entropy-correctness_puzzle.tex, Sec4_stepwise_informativeness_assumption.tex, Sec5_empirical_validation.tex, Sec6_future_work.tex, AppA_experimental_setup.tex, AppB_further_results.tex, AppC_proofs.tex
\documentclass{article}

% Math + basic utilities
\usepackage{amsmath, amsthm, amssymb}
\usepackage{mathtools}
\usepackage{graphicx}
\usepackage{subcaption}

% Tables
\usepackage{array}
\usepackage{tabularx}
\usepackage{makecell}
\usepackage{multicol}
\usepackage{booktabs}

% Hyperlinks
\usepackage{hyperref}

% Colors
\usepackage{xcolor}
\usepackage{xparse}

% --- Theorem environments ---
\newtheorem{assumption}{Assumption}
\newtheorem{theorem}{Theorem}
\newtheorem{remark}{Remark}
\newtheorem{corollary}{Corollary}
\newtheorem{proposition}{Proposition}
\newtheorem{definition}{Definition}
\newtheorem{claim}{Claim}
\newtheorem{lemma}{Lemma}
\newtheorem{hypothesis}{Hypothesis}
\newtheorem{question}{Question}

\usepackage{tcolorbox}
\usepackage{xcolor}

\usepackage[preprint]{icml2026}


% --- New commands ---
\newcommand{\bad}[1]{{\color{failred} #1}}
\newcommand{\weak}[1]{{\color{weakorange} #1}}

% --- Colors ---
\definecolor{summarygray}{RGB}{245,245,245}
\definecolor{summaryborder}{RGB}{80,80,80}
\definecolor{failred}{RGB}{180,30,30}
\definecolor{weakorange}{RGB}{200,120,20}

\newtcolorbox{summarybox}{
  colback=summarygray,
  colframe=summaryborder,
  boxrule=0.6pt,
  arc=3pt,
  left=6pt,
  right=6pt,
  top=6pt,
  bottom=6pt,
  width=\linewidth,
  before skip=8pt,
  after skip=10pt
}

\begin{document}

\icmltitlerunning{Why are Entropy Dynamics and Reasoning Correlated in LLMs?}

\twocolumn[
\icmltitle{The Stepwise Informativeness Assumption: \\ Why are Entropy Dynamics and Reasoning Correlated in LLMs?}

\begin{icmlauthorlist}
\icmlauthor{Mar Gonzàlez I Català}{cam}
\icmlauthor{Haitz Sáez de Ocáriz Borde}{cam,oxf}
\icmlauthor{George D. Monta\~nez}{cam,hmc}
\icmlauthor{Pietro Liò}{cam}
\end{icmlauthorlist}

\icmlaffiliation{cam}{University of Cambridge}
\icmlaffiliation{oxf}{University of Oxford}
\icmlaffiliation{hmc}{Harvey Mudd College}

\icmlcorrespondingauthor{Mar Gonzàlez I Català}{mg2211@cam.ac.uk}

\icmlkeywords{Reasoning, Uncertainty, Entropy, Information Theory}

\vskip 0.3in
]
\printAffiliationsAndNotice{}

\begin{abstract}
Recent work uses entropy-based signals at multiple representation levels to study reasoning in large language models, but the field remains largely empirical. A central unresolved puzzle is why internal entropy dynamics, defined under the predictive distribution of a model, correlate so robustly with external correctness given by the ground-truth answer. In this paper, we argue that this correlation arises because autoregressive models reason correctly when they accumulate information about the true answer via answer-informative prefixes. We formalize this intuition via the Stepwise Informativeness Assumption (SIA), which states that reasoning prefixes accumulate answer-relevant information in expectation as generation progresses. We show that SIA naturally emerges from maximum-likelihood optimization on human reasoning traces and is reinforced by standard fine-tuning and reinforcement-learning pipelines. We then derive observable signatures of SIA linking conditional answer entropy dynamics to correctness. We empirically test SIA across multiple reasoning benchmarks (GSM8K, ARC, SVAMP) and a diverse set of open-weight LLMs (Gemma-2, LLaMA-3.2, Qwen-2.5, DeepSeek and Olmo variants), showing that training induces it and that correct traces exhibit characteristic conditional answer entropy patterns.
\end{abstract}

\vspace{-20pt}
\section{Introduction}

A growing body of empirical studies analyzes model-internal entropy dynamics and consistently reports strong correlations between characteristic patterns and reasoning quality in large language models~(LLMs). These signals have been successfully used to improve reasoning performance \cite{agarwal2025unreasonableeffectivenessentropyminimization, li2025compressingchainofthoughtllmsstep, ton2025understandingchainofthoughtllmsinformation}, guide exploration and early stopping \cite{zhang2025entropybasedexplorationconductionmultistep, sharma2025thinkjustenoughsequencelevel}, identify critical decision points \cite{ali2026entropylensuncoveringdecisionstrategies, wang20258020rulehighentropyminority, qian2025demystifyingreasoningdynamicsmutual}, and detect failures such as hallucinations or overthinking \cite{farquhar_detecting_2024}. However, despite the empirical success of entropy-based approaches to reasoning, a central unresolved question remains:

\begin{question}\label{q:entropy}
Why do internal entropy dynamics—defined purely with respect to a model’s predictive
distribution—correlate so robustly with external correctness, which is defined only
relative to the ground-truth answer?
\end{question}

In this paper, we propose an explanation for this phenomenon. We argue that the observed entropy–correctness correlation arises if autoregressive models learn, through training, to accumulate information about the true answer via answer-informative prefixes, a pattern inherited from human reasoning traces and reinforced by fine-tuning and reinforcement-learning pipelines. We formalize this hypothesis via the \emph{Stepwise Informativeness Assumption} (SIA), a minimal information-theoretic condition stating that reasoning prefixes accumulate information about the true answer in expectation. Under SIA, conditional answer entropy can be interpreted as a progress variable for reasoning: it tracks cumulative answer-relevant information and decreases along successful reasoning chains. Crucially, our framework predicts that characteristic signatures of this descent indicate whether reasoning converges reliably to the correct answer. This provides a structural explanation for why entropy-based signals, despite being internal quantities, can become predictive of reasoning quality.

Finally, we empirically validate the framework across pretrained, supervised fine-tuned, and reinforcement-learning–trained models. We show that (i) training for reasoning induces SIA, and (ii) when SIA holds, it leaves clear traces in entropy dynamics, making conditional answer entropy an informative progress variable.
\section{Preliminaries and Notation}
\label{sec:preliminaries}
We now provide standard definitions of language factorization, LLM training stages, and information theory, on which our results are based.
\subsection{Next-token prediction and likelihood factorization}

Modern language models are trained under the next-token prediction paradigm. 

\begin{definition}[Next-token prediction and autoregressive factorization]
Given an input prefix $X_{1:k}$, a language model with parameters $\theta$ defines a conditional distribution over the next token $p_\theta(X_{k+1}\mid X_{1:k}),$ and the likelihood of a full sequence factorizes autoregressively as $p_\theta(X_{1:K})=\prod_{k=1}^{K} p_\theta(X_k \mid X_{<k})$, where $X_{<k} \coloneqq X_{1:k-1}$.
\end{definition}

\begin{definition}[Autoregressive language model training objective]
Let the training corpus be a collection of $N$ token sequences of variable length $K_i$, $\mathcal{D}=\{X_{1:K_i}^{(i)}\}_{i=1}^N$. The maximum-likelihood training objective for a language model with parameters $\theta$ is defined as $\theta^\ast= \arg \max_\theta \sum^N_{i=1}\log p_\theta(X_{1:K_i}^{(i)})$, which expands autoregressively as a sum over token log-likelihoods.
\end{definition}

In practice, this objective is implemented  by minimizing the cross-entropy loss $\mathcal{L}_{\text{CE}}=-\sum^N_{i=1}\sum^{K_i}_{k=1}\log p_\theta(X_k^{(i)} \mid X_{<k}^{(i)}).$ This encourages the model to make each future token as predictable as possible given the past. Later sections of this paper analyze how this pressure towards subsequent predictability affects reasoning processes, correctness, and entropy minimization.

\subsection{Difference between true answer, chain-of-thought, and model predictive distribution}
\label{subsec:preliminaries-distributions}

The following definitions are key to understanding the difference between the internal model dynamics and the ground-truth distribution referenced in Question~\ref{q:entropy}.

\begin{definition}[True answer distribution] Let $Q \in \mathcal{Q}$ denote a query and $A\in \mathcal{A}$ its correct answer. The ground-truth joint distribution over queries and answers is $(Q,A)\sim p^\star(Q,A)$, and the corresponding true posterior over answers given a query is $p^\star(A \mid Q)$. All statements about correctness are defined with respect to this true conditional distribution $p^\star(A \mid Q)$.
\label{def:trueanswer}
\end{definition}

\begin{definition}[Chain-of-thought (data-generating) distribution] In many reasoning datasets, each query $Q$ is paired with a correct answer $A$ and a human-written chain-of-thought trace $C_{1:K}$. We denote the empirical joint distribution over this triple as $r(Q,C_{1:K},A) = p^\star(Q,A)\, r(C_{1:K}, A\mid Q)$,
where $p^\star(Q,A)$ is the ground-truth question--answer distribution and  $r(C_{1:K}, A\mid Q)$ describes how human annotators produce chain-of-thought traces when solving the problem.
\label{def:datadistribution}
\end{definition}

\begin{definition}[Model predictive distribution] Given a query $Q$, a reasoning model with parameters $\theta$ generates a sequence of intermediate tokens $C_{1:K} = (C_1, \dots, C_K)$ and an answer sequence $A= (A_1,\dots,A_T)$. The model induces an autoregressive distribution over full reasoning traces, $p_\theta(C_{1:K} \mid Q) = \prod_{k=1}^K p_\theta(C_k \mid Q, C_{<k})$, and, conditioned on a reasoning trace, an autoregressive distribution over answers, $p_\theta(A \mid Q, C_{1:K})=\prod^T_{t=1} p_\theta(A_t \mid Q, C_{1:K}, A_{<t})$.
\label{def:modelpredictive}
\end{definition}

Note that we abuse notation by using $A$ to denote both the model-generated and ground-truth answer. The intended meaning will be clear from the underlying distribution.

When defining stepwise entropy and information-gain quantities, we will also condition on partial prefixes $C_{1:k}$ (for $k<K$), which yields $p_\theta(A\mid Q,C_{1:k})$ by the same factorization. Importantly, note that token-level entropies and conditional answer entropies are purely internal properties of the model's internal predictive distribution $p_\theta$. These entropies are in principle independent of the true external answer distribution $p^\star(A\mid Q)$.



\subsection{Training stages of language models}

InstructGPT \cite{ouyang2022traininglanguagemodelsfollow} formalized a three-stage training pipeline that has since become standard in modern language models.

\paragraph{Pretraining on raw data.} In the pretraining stage, the model is trained via maximum-likelihood estimation on large-scale text corpora using the next-token prediction objective previously described. This implicitly includes a wide variety of reasoning traces such as explanations, derivations, proofs, and step-by-step problem solutions. Although correctness is not explicitly optimized at this stage, the model is rewarded for generating continuations that make future tokens predictable given the past, thereby learning sequential structures that progressively constrain plausible outcomes.

\paragraph{Supervised fine-tuning on labeled chain-of-thought triples.} In supervised fine-tuning (SFT), the model is trained on datasets consisting of explicit triples $(Q, C_{1:K}, A)$, where $C_{1:K}$ is a human-written chain-of-thought leading to the correct answer $A$. The same maximum-likelihood objective is applied, but now correctness is directly reflected in the data distribution: reasoning traces that make the correct answer highly probable receive higher likelihood. As a result, the model is explicitly encouraged to generate intermediate steps that reduce uncertainty about the true answer.

\paragraph{Post-training with reinforcement learning.} Finally, it is common to apply reinforcement learning–based post-training methods such as PPO~\citep{schulman2017proximal}, GRPO~\citep{shao2024deepseekmath}, or RL with verifiable rewards~(RLVR)~\citep{wen2025reinforcementlearningverifiablerewards} to elicit reasoning in LLMs. These methods reweight or refine the model’s generation policy based on outcome-level or process-level reward signals, further reinforcing reasoning trajectories that lead to correct answers and penalizing those that do not. This stage strengthens the alignment between internal uncertainty reduction and external correctness, but does not introduce new reasoning primitives; rather, it reshapes the probability mass over existing reasoning patterns learned during pretraining and SFT.

\subsection{Information-Theoretic Preliminaries}

When an LLM reasons step by step, each intermediate token can raise or lower its confidence in the correct answer. Information theory provides a principled framework to quantify these changes, formalizing uncertainty and information gain in probabilistic systems.

Next, we summarize the information-theoretic measures used throughout the paper. All random variables are assumed to be discrete. Also, before introducing the definitions, we clarify our notation: uppercase letters (e.g., $A$, $Q$, $C_k$) denote random variables; lowercase letters (e.g., $a$, $q$, $c_k$) denote particular realizations or sampled values of those variables; calligraphic letters (e.g., $\mathcal{A}$, $\mathcal{Q}$, $\mathcal{C}$) denote the set of all possible values each random variable can take. Unless otherwise stated, all logarithms are natural logarithms.

\begin{definition}[Entropy]
    Let $X$ be a discrete random variable taking values in $\mathcal{X}$, with probability mass function $p(x)=\Pr[X=x]$. The entropy (average or expected surprisal) of $X$ is defined as $H(X):=-\sum_{x\in\mathcal{X}} p(x)\log p(x).$
\end{definition}

\begin{definition}[Conditional entropy]
    Let $X$ and $Y$ be discrete random variables taking values in $\mathcal{X}$ and $\mathcal{Y}$, with joint pmf $p(x,y)=\Pr[X=x,Y=y]$, and marginal pmfs $p(x)=\Pr[X=x]$, $p(y)=\Pr[Y=y]$. The \emph{conditional entropy} of $Y$ given $X$ is $H(Y\mid X):=-\sum_{x\in\mathcal{X},y\in\mathcal{Y}}p(x,y)\log \left(\frac{p(x,y)}{p(x)}\right),$
    with the convention that $0\log0=0$.
\end{definition}

\begin{definition}[Mutual Information]
    Let $X$ and $Y$ be discrete random variables taking values in $\mathcal{X}$ and $\mathcal{Y}$, with joint pmf $p(x,y)=\Pr[X=x,Y=y]$, and marginal pmfs $p(x)=\Pr[X=x]$, $p(y)=\Pr[Y=y]$. The \emph{mutual information between $X$ and $Y$} is defined as $I(X;Y):=\sum_{x\in\mathcal{X}, y\in \mathcal{Y}}p(x,y)\log\left(\frac{p(x,y)}{p(x)p(y)}\right).$
\end{definition}

Note that all these definitions rely on logarithms. While the use of logarithms is not mandated, they uniquely satisfy a few intuitive properties: information from independent events adds, rarer events carry more information, and small changes in probability produce small changes in information.

\section{Why does entropy track correctness in reasoning models?}
\label{sec:entropy-correctness_puzzle}

Internal entropy is defined entirely under a model’s predictive distribution (Definition~\ref{def:modelpredictive}), whereas correctness is defined with respect to an external ground-truth answer distribution (Definition~\ref{def:trueanswer}). There is therefore no \textit{a priori} reason for these two notions to be aligned: internal uncertainty could track stylistic variability, spurious hypotheses, or model-internal ambiguity unrelated to task success. Indeed, recent work cautions against treating intermediate tokens as faithful indicators of reasoning progress or task difficulty \cite{kambhampati2025stopanthropomorphizingintermediatetokens, palod2025performativethinkingbrittlecorrelation}.

\subsection{Empirical evidence for entropy-correctness alignment}

Despite this conceptual gap, numerous studies report a robust correlation between internal entropy dynamics and reasoning accuracy, across tasks, model families, and levels of granularity. This correlation is exploited for analysis, control, and prediction of reasoning behavior.

\paragraph{Analysis.}
High entropy is associated with overextrapolation (``hallucination'') and unreliable outputs, while entropy plateaus correspond to ``overthinking,'' where additional reasoning does not improve accuracy \cite{farquhar_detecting_2024}. Successful trajectories exhibit distinctive entropy patterns: uncertainty concentrates at critical “forking” steps and is systematically reduced thereafter \cite{qian2025demystifyingreasoningdynamicsmutual, wang20258020rulehighentropyminority}.

\paragraph{Control.}
Early-stopping methods terminate chain-of-thought generation once entropy plateaus or falls below a threshold \cite{ sharma2025thinkjustenoughsequencelevel}, while compression and exploration-based approaches treat entropy as a signal of decision points, pruning or expanding reasoning accordingly \cite{li2025compressingchainofthoughtllmsstep, zhang2025entropybasedexplorationconductionmultistep}.

\paragraph{Prediction.}
Entropy-based metrics can reliably predict whether an ongoing reasoning trajectory will ultimately be correct: traces with a decreasing entropy trajectory are much more likely to end in correct answers \cite{guo2025measuringreasoningutilityllms, liu2025tokensignaturepredictingchainofthought}.

Thus, internal uncertainty dynamics track external correctness closely enough that many methods implicitly treat entropy reduction as a proxy for reasoning progress. But why should this be true at all?

\subsection{Common justifications for entropy-based reasoning methods}

The literature offers several recurring explanations, none of which fully resolve the puzzle. A common implicit assumption is that reductions in entropy reflect a narrowing of the space of plausible solutions \cite{qian2025demystifyingreasoningdynamicsmutual, ton2025understandingchainofthoughtllmsinformation}. This interpretation presupposes that the uncertainty being reduced concerns the correct answer.

Other works appeal to training-induced alignment: since models are trained to produce correct answers, their internal uncertainty should track correctness \cite{sharma2025thinkjustenoughsequencelevel}. This would be compelling if it specified the structural properties of the learned distribution that ensure predictive entropy becomes aligned with the ground truth throughout a reasoning chain, but such conditions are not articulated.

Some analyses assume that reasoning steps reduce uncertainty about the true hypothesis \cite{ton2025understandingchainofthoughtllmsinformation}. However, this presupposes a coupling between intermediate model states and the ground-truth answer that is not derived from the training objective or the structure of the learned distribution.

Finally, many works offer no justification at all, treating the entropy--correctness correlation as an empirical fact to be exploited rather than a phenomenon to be explained \cite{liu2025tokensignaturepredictingchainofthought, zhang2025entropybasedexplorationconductionmultistep}. Exploiting a correlation, however, does not explain it. To our knowledge, no prior work asks why this alignment should arise, or under what conditions it should be expected to hold or fail.



\section{Stepwise Informativeness Assumption}
\label{sec:assumption}

To explain when internal uncertainty reflects external correctness, we formalize a minimal mechanism by which reasoning prefixes come to encode information about the true answer. Proofs for all lemmata, propositions, and theorems can be found in Appendix~\ref{app:proofs}.

\subsection{Stepwise information gain}
We introduce local, token-level quantities that capture how individual reasoning steps affect uncertainty about the answer. These quantities allow us to describe reasoning progress at the granularity of single tokens, before aggregating to prefix-level information.

\begin{definition}[Pointwise surprisal]
    For a sampled triple $(Q=q,A=a,C_{1:k}=c_{1:k})$, we define the \emph{pointwise conditional surprisal} as:
    $h(a\mid q,c_{<k})=-\log p(a \mid q, c_{<k})$
\end{definition}

\begin{definition}[Information gain]
    For a sampled triple $(Q=q,A=a,C_{1:k}=c_{1:k})$, we define the  \emph{pointwise information gain of step $k$} as $
    \Delta_k(q,a,c_{1:k}):=h(a\mid q,c_{<k})-h(a\mid q,c_{\leq k})
    \label{eq:pointwise-delta}$
\end{definition}

\begin{remark}
The interpretation of this quantity is: $\Delta_k>0$, the step makes the correct answer more probable (informative step); $\Delta_k<0$, the step makes the correct answer less probable (misinformative step).
\end{remark}

\begin{lemma}
The expected value of $\Delta_k$ equals the standard conditional mutual information: $\mathbb{E}[\Delta_k] = I(A; C_k \mid Q, C_{<k}) = H(A \mid Q, C_{<k}) - H(A \mid Q, C_{\le k}).$
\label{cl:exp_delta_is_I}
\end{lemma}

\begin{definition}[Cumulative gain]
    For a sampled triple $(Q=q,A=a,C_{1:k}=c_{1:k})$, we define the \emph{cumulative gain up to step $k$} as $G_k := \sum_{t=1}^k \Delta_t=h(a\mid q)-h(a\mid q,c_{\leq k}).$ In expectation, $\mathbb{E}[G_k]=I(A;C_{1:k} \mid Q)=\sum_{t=1}^k I(A;C_t \mid Q, C_{<t}).$
    \label{def:cumulativegain}
\end{definition}

Lastly, note that entropy and mutual-information quantities are always understood with respect to an underlying probability distribution. To make this explicit, we will attach the distribution as a subscript whenever there is ambiguity, e.g. $H_{p}(\cdot)$ and $I_{p}(\cdot)$ for a given probability distribution $p$. Also, we assume stochastic decoding from $p_\theta$. (Deterministic) greedy decoding is a degenerate case, since it selects the most probable token at each step, often collapsing token-level entropy and trivializing many of the entropy-based quantities studied, which obscures stepwise uncertainty dynamics.

\subsection{Stepwise Informativeness Assumption}

To relate model-internal uncertainty to external correctness, we introduce a joint coupling between the model’s reasoning traces and the ground-truth answer distribution. We consider
\(
\Pi := \{p(Q,C_{1:K},A) : p(Q,A)=p^\star(Q,A),\; p(C_{1:K}\mid Q)=p_\theta(C_{1:K}\mid Q)\},
\)
where $p^\star$ denotes the ground-truth query–answer distribution and $p_\theta$ the model’s predictive distribution over reasoning traces. This avoids imposing any conditional independence between $A$ and $C_{1:K}$ given $Q$.

\begin{proposition}[Conditional answer entropy as cumulative information]
\label{prop:entropy-information}
Under a fixed joint $p\in \Pi$ and for any prefix length $k \ge 1$,
\(
H_p(A \mid Q, C_{1:k})
=
H_p(A \mid Q) - \sum_{t=1}^k I_p(A; C_t \mid Q, C_{<t})
=
H_p(A \mid Q) - I_p(A; C_{1:k} \mid Q)
\)
\end{proposition}

Thus, under $p$, conditional answer entropy is not merely an internal uncertainty measure: it is a progress variable tracking how much information about the true answer has been accumulated.

This motivates a structural assumption linking entropy dynamics to correctness.

\begin{assumption}[Stepwise Informativeness Assumption (SIA)]
\label{ass:SIA}
Under a fixed joint $p\in\Pi$, prefixes are informative about the true answer in expectation:
\[
I_p(A;C_{1:k}\mid Q)
\ge \epsilon_k > 0
\quad
\text{for all } k\ge 1.
\] 
\end{assumption}

We call this the Stepwise Informativeness Assumption (SIA) because it implies that partial reasoning prefixes contain information about the final answer. Note that SIA does not require each individual reasoning token to be informative; rather, it constrains the total information contained in the prefix $C_{1:k}$. This formulation accounts for redundant tokens or stalling phrases that do not provide immediate marginal information but are part of a larger informative prefix.

SIA is a property of the joint coupling $p$, not of $p_\theta$ alone. Entropy reduction under the model’s internal posterior does not imply SIA unless prefixes are also informative about the true answer under an answer-consistent coupling. In the absence of SIA, conditional entropy may decrease for purely internal reasons while correctness does not improve.

When $\exists p \in \Pi$ and SIA holds, entropy-based reasoning diagnostics are theoretically justified: the sequence $\{\epsilon_k\}$ quantifies cumulative answer-relevant information gain, and the trajectory $H_p(A\mid Q,C_{1:k})$ characterizes whether the model is progressing toward the correct answer.


\subsubsection{Entropy constrains achievable accuracy}
\label{subsec:entropy_constrains}

\begin{theorem}[Entropy constrains achievable accuracy]
\label{th:entropy_accuracy}
Under a fixed joint $p\in \Pi$ and for any prefix length $k \ge 1$, let $\widehat A_k$ denote the Bayes-optimal predictor
based on $(Q, C_{1:k})$ under the posterior $p(A \mid Q, C_{1:k})$, and let $P_e^{(k)} := \Pr(\widehat A_k \neq A)$ denote its misclassification probability. Then
\[
P_e^{(k)}
\;\ge\;
\frac{H_p(A \mid Q, C_{1:k}) - \log 2}{\log(|\mathcal{A}| - 1)}, \qquad \text{where} |\mathcal{A}|>2
\]
\end{theorem}


This bound shows that correctness is limited by how informative reasoning prefixes are about the true answer: prefixes that substantially reduce conditional answer entropy yield lower error, while weakly informative prefixes cannot support high accuracy, regardless of the predictor.

Theorem~\ref{th:entropy_accuracy} gives a necessary condition for correctness: a reasoning chain cannot be reliably correct unless its prefixes exhibit sufficiently low conditional answer entropy.

\subsubsection{Early vs.\ late information gain}
\label{subsec:early_vs_late}

Consider two reasoning chains that satisfy SIA. If one chain attains lower conditional answer entropy than the other over an initial segment of the reasoning trace, then throughout that segment it admits a strictly lower information-theoretic lower bound on achievable error (Theorem~\ref{th:entropy_accuracy}). Even when total information gain is matched by the end of the trace, earlier entropy reduction yields a larger fraction of tokens generated under low conditional entropy, where downstream steps are less likely to be derailed by sampling noise or spurious branches.

This leads to an operational criterion for detecting correct chains: correct reasoning chains should “lock onto” the answer early, before they are forced to by the monotonicity of conditional entropy.

\subsubsection{Saturation}

For many tasks, the total amount of answer-relevant information that can be extracted from a reasoning trace is finite. As conditional answer entropy decreases and approaches its minimum, the amount of remaining answer-relevant uncertainty necessarily shrinks. Consequently, any further reductions in conditional entropy must become progressively smaller and may eventually be negligible. When this occurs, conditional entropy effectively plateaus: additional reasoning steps cannot meaningfully reduce uncertainty about the answer. 

Reaching a plateau is not sufficient for correctness (as incorrect chains may also saturate around an erroneous hypothesis), but failure to saturate constitutes negative evidence against correctness.

\subsection{Why is SIA a reasonable assumption?}
\label{subsec:SIA_reasonable}

SIA is not guaranteed to hold universally: prefixes might not be informative about the true answer. But why is it reasonable to expect it to hold?

\subsubsection{Stepwise Informativeness in human-generated reasoning traces}

Human reasoning traces often exhibit progressive accumulation of answer-relevant information, even without explicit optimization for correctness. This follows from general constraints on sequential information processing.

Recent research  \cite{futrell_linguistic_2025} shows that under realistic cognitive constraints (limited memory, attention, and processing capacity) sequential signals that minimize \textit{predictive information}, the mutual information between past and future, adopt a characteristic structure: information is decomposed into approximately independent components that are expressed locally and incrementally. This yields sequences that are systematically and progressively informative, closely matching the structure of natural language and assisting in downstream sequence prediction. 

Human-written reasoning traces are a special case of such sequential signals, with the additional property that their future includes the correct answer. Under the same constraints, earlier prefixes are therefore expected to increasingly constrain the space of plausible continuations and answers. As reasoning unfolds, the correct answer becomes more predictable in aggregate.

Formally, let \(C_{1:K}\) denote a human-generated chain-of-thought and \(A\) the correct answer. If reasoning traces minimize predictive information, then prefixes \(C_{1:k}\) carry increasing mutual information about future tokens, including \(A\). Equivalently, under a data generating distribution $r(Q,C_{1:K},A)$, the conditional answer entropy \(H_r(A \mid Q, C_{1:k})\) decreases with \(k\) in expectation, implying growing prefix-level mutual information \(I_r(A; C_{1:k} \mid Q)\).

Crucially, this argument does not assume that humans optimize intermediate steps for correctness or have access to the answer distribution during generation. Stepwise informativeness instead emerges as a structural consequence of general cognitive pressures on sequential communication.

\subsubsection{Transfer of Stepwise Informativeness under Maximum Likelihood Training}

We now study whether stepwise informativeness present in human-generated reasoning traces transfers to a model via MLE training. 

\begin{lemma}
\label{lemma:log-likelihood_identity}
Let $r$ denote the empirical data-generating distribution over full sequences 
$X = (Q,C_{1:K},A)$, and let $p_\theta$ be the model distribution.  
The negative log-likelihood objective $\mathcal{L}(\theta) \;=\; \mathbb{E}_{X\sim r}\!\left[-\log p_\theta(X)\right]$
satisfies the identity
\[
\mathcal{L}(\theta)
= -\sum_{x} r(x)\log p_\theta(x)
= H_r(\cdot) + \mathrm{KL}\!\left(r \,\Vert\, p_\theta\right),
\]
where $H_r(\cdot)$ is independent of $\theta$.  
Hence minimizing $\mathcal{L}(\theta)$ is equivalent to minimizing 
$\mathrm{KL}(r \,\Vert\, p_\theta)$, and thus drives $p_\theta$ toward $r$ in forward KL-divergence, up to model capacity.
\end{lemma}

\begin{lemma}[KL Decomposition of the Joint Conditional]
\label{lemma:KLDecompositionoftheJointConditional}
For any joint distributions $r$ and $p_\theta$ over $(C_{1:K},A)$ conditioned on $Q$,
the following identity holds:
\begin{equation*}
\begin{aligned}
&\mathrm{KL}\!\left(
r(C_{1:K},A\mid Q)
\;\Vert\;
p_\theta(C_{1:K},A\mid Q)
\right)
= \\
&= \mathrm{KL}\!\left(r(C_{1:K}\mid Q)\,\Vert\,p_\theta(C_{1:K}\mid Q)\right)
+ \\
&\mathbb{E}_{r(C_{1:K}\mid Q)}
\!\left[
\mathrm{KL}\!\left(
r(A\mid Q,C_{1:K})
\,\Vert\,
p_\theta(A\mid Q,C_{1:K})
\right)
\right].
\end{aligned}
\end{equation*}
\label{lemma:joint-kl-decomposition}
\end{lemma}

\begin{lemma}[MLE Implies Marginal and Conditional Alignment]
\label{lemma:MLE_implies_alignment}
Let $Q$ be any question in the support of $r(Q)$.
Given $\delta>0$, if
\(
\mathrm{KL}\!\left(
r(C_{1:K},A \mid Q)\,
\Vert\,
p_\theta(C_{1:K},A \mid Q)
\right)
\le \delta,
\label{eq:joint-smallness}
\)
then 
\(
\mathrm{KL}\!\left(
    r(C_{1:K}\mid Q)
    \,\Vert\,
    p_\theta(C_{1:K}\mid Q)
    \right)
    \le \delta
\) and 
\(
\mathbb{E}_{r(C_{1:K}\mid Q)}
    \!\left[
    \mathrm{KL}\!\left(
    r(A\mid Q,C_{1:K})
    \,\Vert\,
    p_\theta(A\mid Q,C_{1:K})
    \right)
    \right]
    \le \delta .
\)
\end{lemma}

The formal guarantee relies on continuity of entropy and conditional mutual information on finite alphabets. To relate informativeness under the data-generating distribution $r$ to that under the model distribution $p_\theta$, we require that entropy be stable under small distributional perturbations. 

\begin{lemma}[Continuity of Entropy under KL]\label{lemma:entropy-continuity}
Let $P$ and $Q$ be probability distributions on a finite alphabet
$\mathcal{X}$ satisfying $\mathrm{KL}(P \Vert Q) \le \delta.$
Then there exists a function $f_\mathcal{X} : [0,\infty) \to [0,\infty)$
with $f_\mathcal{X}(\delta) \to 0$ as $\delta \to 0$ such that $\bigl| H(P) - H(Q) \bigr| \le f_\mathcal{X}(\delta).$ In particular, for every $\varepsilon > 0$ there exists
$\delta > 0$ such that $\mathrm{KL}(P \Vert Q) \le \delta \quad\Longrightarrow\quad
\bigl| H(P) - H(Q) \bigr| \le \varepsilon.$
\end{lemma}

The same argument extends to conditional entropy by applying continuity to the relevant joint and marginal distributions. 

\begin{lemma}[Continuity of Conditional Entropy]\label{lemma:cond-entropy-continuity}
Let $P$ and $Q$ be distributions on a finite product alphabet
$\mathcal{X} \times \mathcal{Y}$, with $\mathrm{KL}(P \Vert Q) \le \delta.$
Then there exists a function $g_{\mathcal{X},\mathcal{Y}}(\delta)$
with $g_{\mathcal{X},\mathcal{Y}}(\delta) \to 0$ as $\delta \to 0$
such that $\bigl| H_P(Y \mid X) - H_Q(Y \mid X) \bigr|
\le g_{\mathcal{X},\mathcal{Y}}(\delta).$
Equivalently, for every $\varepsilon > 0$ there exists $\delta > 0$
such that $\mathrm{KL}(P \Vert Q) \le \delta
\quad\Longrightarrow\quad
\bigl| H_P(Y \mid X) - H_Q(Y \mid X) \bigr| \le \varepsilon.$
\end{lemma}

Combining these results yields continuity of conditional mutual information, which enables internal stepwise informativeness transfer from $r$ to $p_\theta$ up to an arbitrarily small error.

\begin{lemma}[Continuity of Conditional Mutual Information]\label{lemma:mi-continuity}
Let $r$ and $p_\theta$ be distributions on a finite product alphabet
$\mathcal{Q}\times\mathcal{C}_1\times\cdots\times\mathcal{C}_K
\times\mathcal{A}$, and fix $k\in\{1,\dots,K\}$. Suppose that $\mathrm{KL}(r \Vert p_\theta) \le \delta.$ Then there exists a function $G_k(\delta)$ with $G_k(\delta)\to 0$ as
$\delta\to 0$ such that $\bigl|
I_r(A; C_{\le k} \mid Q)
-
I_{p_\theta}(A; C_{\le k} \mid Q)
\bigr|
\le G_k(\delta),$ where the mutual informations and entropies on the right-hand side are computed under $p_\theta$, and those on the left-hand side under $r$. Equivalently, for every $\varepsilon>0$ there exists $\delta>0$ such
that $\mathrm{KL}(r \Vert p_\theta) \le \delta
\quad\Longrightarrow\quad
\bigl|
I_r(A; C_{\le k} \mid Q)
-
I_{p_\theta}(A; C_{\le k} \mid Q)
\bigr|
\le \varepsilon.$
\end{lemma}


\begin{theorem}[Transfer of internal stepwise informativeness to the model]
\label{thm:transfer-sia}
Given a step $k \geq 1$, let $r$ denote the empirical data joint over $(Q,C_{1:K},A)$, and suppose that $I_r(A;C_{\leq k} | Q)>0$. Let $p_\theta$ denote the model distribution over $(Q,C_{1:K},A)$, then there exists $\delta_k > 0$ such that, whenever
\(
\mathrm{KL}(r \,\Vert\, p_\theta) \le \delta_k,
\)
the model joint $p_\theta$ satisfies
\(
I_{p_\theta}(A; C_{\leq k} \mid Q) \ge \frac{\varepsilon_k}{2} > 0.
\)
\end{theorem}

The transfer result establishes that if the data-generating distribution $r$ exhibits stepwise informativeness, then a model trained under MLE will inherit an \emph{internal} version of this property, which does not by itself imply SIA.

Nonetheless, when supervision consists of explicit triples $(Q,C_{1:K},A)$, the objective has a well-defined target: the correct answer $A$. Under MLE, prefixes that systematically increase the probability of $A$ are reinforced, and predictive information is therefore concentrated on intermediate steps that progressively constrain the answer space toward correctness. In contrast, during large-scale pretraining, reasoning-like continuations are embedded in a corpus where next-token prediction is governed by distributional regularities rather than any particular ground-truth objective. As a result, the model may learn to produce locally coherent reasoning patterns without those prefixes being systematically informative about a true answer variable. Thus, while both regimes optimize next-token predictability, only SFT systematically ties predictive information to answer-relevant structure, making SIA behavior empirically more likely after supervision than after pretraining alone. 

\subsection{Regimes in which SIA does not hold}
\label{subsec:sia_not_hold}

Entropy-based diagnostics are not theoretically justified if training fails to induce an answer-compatible distribution $p \in \Pi$ that satisfies SIA and that $p_\theta$ faithfully approximates.

In this case, conditional answer entropy under $p_\theta$ may decrease along a reasoning trace even as the model converges to an incorrect answer. Formally, such trajectories satisfy an \emph{internal stepwise informativeness} condition:
\(
I_{p_\theta}(A ; C_{\leq k} \mid Q) > 0,
\)
despite vanishing informativeness under the joint distribution $p$ induced by training. Entropy descent then reflects uncertainty reduction with respect to a misaligned belief state, providing an information-theoretic formalization of ``hallucinations'', common in adversarial, out-of-distribution, or weakly supervised settings. 

Lastly, it is worth noting the theory behind SIA is most applicable to problems with a well-defined terminal variable, such as mathematical reasoning or multiple-choice question answering, as opposed to free-form outputs like creative writing.

\begin{summarybox}
\textbf{Summary}: we have: (i) introduced SIA as a minimal, falsifiable condition and structural theory under which entropy-based reasoning analyses are justified, (ii) shown that MLE induces an internal form of SIA, (iii) used KL-continuity to justify transfer, and (iv) explicitly studied what training does not guarantee.
\end{summarybox}
\section{Empirical validation}

In this section, we test whether training induces an answer-consistent $p\in\Pi$ that $p_\theta$ faithfully approximates, and under what conditions. Empirically, we do not directly verify SIA, which is a property of the joint coupling $p \in \Pi$. Instead, we evaluate the entropy dynamics predicted by SIA and ask whether training induces model behavior compatible with such a coupling.

We organize our empirical evaluation around three questions: (i) does conditional entropy descent align with increasing probability of the true answer, (ii) is this alignment induced and strengthened by training for reasoning, and (iii) what are the observable signatures and failure modes of SIA?

We evaluate eleven models across three datasets (GSM8K, ARC, SVAMP), spanning base, instruction-tuned, CoT-tuned and RL-trained regimes. All entropy quantities are estimated via Monte Carlo rollouts under stochastic decoding. Full evaluation details are provided in Appendix~\ref{app:exp}.

\subsection{Entropy-answer alignment}

If training has successfully aligned the model’s internal joint with a coupling $p\in\Pi$ that satisfies SIA, reductions in conditional answer entropy should coincide with increases in the probability assigned to the true answer. To test this directly, we define the following diagnostic.

\paragraph{SIA alignment coefficient.}
For each generated trace, we compute the correlation (across prefix steps~$k$) between conditional answer entropy and gold surprisal:
\begin{equation*}
\rho_{\mathrm{SIA}} \;:=\; \mathrm{corr}_{k}\!\Big(
H_{\theta}(A \mid q, c_{\le k}),
\;-\log p_{\theta}(a^{\star} \mid q, c_{\le k})
\Big).
\end{equation*}
Positive $\rho_{\mathrm{SIA}}$ indicates that uncertainty reduction is aligned with increasing probability of the correct answer, suggesting the internal entropy descent is compatible with an answer-consistent coupling in $\Pi$ that satisfies SIA. Negative values indicate confident misalignment: entropy decreases while moving away from the true answer.

Table~\ref{tab:sia_alignment_by_dataset} summarizes $\rho_{\mathrm{SIA}}$ by model. Base models frequently exhibit weak or negative alignment, whereas supervised fine-tuned models show strong positive alignment on average and RL-trained models approach near-perfect alignment. This indicates that truth-directed entropy descent is not a generic property of autoregressive models, but a training-induced structural feature.

Within each training stage, alignment varies with data curation and optimization objectives. Among base models, Qwen2.5-3B exhibits stronger alignment than Gemma-2 and LLaMA-3.2, probably due to a pretraining corpus richer in reasoning text. Within SFT models, DeepSeek-Chat underperforms, which may be caused by supervision that prioritizes conversational helpfulness. Finally, models explicitly optimized for reasoning, such as OLMo and DeepSeek-R1, exhibit near-perfect alignment, reflecting training regimes that strongly couple intermediate steps to the correct answer.

\begin{table}[hbtp!]
\caption{\textbf{Training aligns entropy descent with the true answer.}
We report the correlation between conditional answer entropy and gold surprisal along each trace
(SIA alignment coefficient $\rho_{\mathrm{SIA}}$), averaged by model and dataset.
\bad{Negative} or \weak{near-zero} values indicate failure of alignment.}
\centering
\small
\setlength{\tabcolsep}{4pt}
\begin{tabular}{lcccc}
\toprule
\textbf{Model} & \textbf{Training} &
\textbf{GSM8K} & \textbf{SVAMP} & \textbf{ARC} \\
\midrule
Qwen2.5-3B        & Base & 0.682 & 0.603 & 0.344 \\
Qwen2.5-3B-it     & SFT  & 0.744 & 0.835 & 0.666 \\
Qwen2.5-Math-1.5B & SFT  & 0.499 & 0.802 & 0.676 \\
DeepSeek-Chat-7B  & SFT  & 0.346 & 0.295 & 0.143 \\
DeepSeek-R1-Distilled & SFT+RL & 0.795 & 0.593 & 0.783 \\
Gemma-2-2B        & Base & \bad{-0.530} & 0.169 & \bad{-0.208} \\
Gemma-2-2B-it     & SFT  & 0.522 & 0.462 & 0.578 \\
LLaMA-3.2-3B      & Base & \bad{-0.361} & 0.424 & \bad{-0.366} \\
LLaMA-3.2-3B-it   & SFT  & 0.576 & 0.399 & 0.545 \\
Olmo-3-7B-Think-SFT & SFT & 0.964 & 0.884 & 0.960 \\
Olmo-3-7B-Think   & SFT+RL & 0.885 & 0.778 & 0.887 \\
\bottomrule
\end{tabular}
\label{tab:sia_alignment_by_dataset}
\end{table}

\subsection{Observable signatures: early lock-in, separability, and saturation}
\label{sec:signatures}

When training has successfully internalized SIA, it gives rise to observable token-level signatures (as often reported in the literature) that distinguish aligned from non-aligned models, and correct from incorrect traces.

\paragraph{Early information accumulation.}
Figure~\ref{fig:gain_correct_incorrect} plots a normalized version of the cumulative information gain (Definition~\ref{def:cumulativegain}),
\(
G(s):=\frac{I(A;C_{\le s}\mid Q)}{I(A;C_{\le 1}\mid Q)}
\)
split by correctness. Correct traces accumulate a larger fraction of their total answer-relevant information earlier in the generation. As predicted by Theorem~ \ref{th:entropy_accuracy}, prefixes with lower entropy are more likely to lead to correct answers. This signature is not observed in non-aligned models (see Appendix~\ref{subsec:signatures_weaken}).

\begin{figure}[hbtp!]
\centering
\includegraphics[width=0.75\linewidth]{figures/normalized_gain_curve.png}
\caption{\textbf{Early information accumulation.}
Normalized cumulative gain $G(s)$ vs.\ relative prefix length $s$, split by correctness in llama-3.2-3B-it (aligned model) in GSM8k dataset.}
\label{fig:gain_correct_incorrect}
\end{figure}

\paragraph{Early separability of correct vs.\ incorrect traces.}
Figure~\ref{fig:auc_by_r} reports the AUC for using conditional entropy at prefix length $s$ to distinguish correct from incorrect traces. For SIA-internalized models, separability is already strong well before the answer is produced, showing that entropy becomes diagnostic early in the trace. This signature is not observed in non-aligned models (see Appendix~\ref{subsec:signatures_weaken}).

\begin{figure}[hbtp!]
\centering
\includegraphics[width=0.75\linewidth]{figures/separability_auc_by_r.png}
\caption{\textbf{Separability.}
AUC for using conditional answer entropy to distinguish correct from incorrect traces vs.\ relative prefix length $s$ across aligned models in GSM8k dataset.}
\label{fig:auc_by_r}
\end{figure}

\paragraph{Saturation.}
Finally, Figure~\ref{fig:entropy_plateaus} shows mean entropy trajectories across model families. Aligned models reach plateaus at (near-)zero conditional answer entropy, consistent with exhausting answer-relevant information, while non-aligned models stabilize at nonzero entropy and exhibit late-stage rebounds, indicating that uncertainty ceases to decrease without converging to a specific answer.

\begin{figure}[hbtp!]
\centering
\includegraphics[width=0.75\linewidth]{figures/entropy_curves.png}
\caption{\textbf{Saturation.}
Mean conditional answer entropy trajectories across non-aligned and aligned models in GSM8k dataset.}
\label{fig:entropy_plateaus}
\end{figure}


Together, these patterns characterize SIA-internalized reasoning: entropy both constrains achievable accuracy and reveals when and how answer-relevant information is acquired. Importantly, all signatures vanish or weaken when this structure is absent (see Appendix~\ref{subsec:signatures_weaken}). 

\subsection{Ablations}
\label{sec:ablations}

Finally, we test whether observed dynamics reflect stepwise structure rather than superficial artifacts.

\paragraph{Shuffle-prefix ablation (post-hoc).}
Table~\ref{fig:shuffle_ablation} shows that randomly permuting tokens within prefixes (length preserved) sharply degrades alignment, indicating that truth-directed entropy descent depends on structured accumulation rather than token count. This permutation is applied only at evaluation time when computing conditional answer distributions and the associated entropies, and does not affect generation.

\begin{table}[hbtp!]
\caption{\textbf{Shuffle-prefix ablation}. Entropy--correctness alignment ($\rho_{\text{SIA}}$) drops sharply when prefix tokens are permuted. \bad{Negative} or \weak{near-zero} values indicate coupling misalignment.}
\label{fig:shuffle_ablation}
\centering
\small
\begin{tabular}{lcc}
\toprule
\textbf{Model} & {Original mean} & {Shuffled mean} \\
\midrule
Qwen2.5-3B      &  0.682 & \bad{-0.132} \\
Qwen2.5-3B-it   &  0.744 & \bad{-0.005} \\
DeepSeek-R1-Distilled    &  0.795 &  \weak{0.020} \\
Gemma-2-2B-it   &  0.522 & \bad{-0.063} \\
\bottomrule
\end{tabular}
\end{table}

Further ablations can be found in Appendix~\ref{subsec:further_ablations}.

\section{Conclusion and Open Questions}
\label{sec:open_quest}

This work provides a structural explanation for why internal entropy dynamics correlate with correctness in autoregressive reasoning models. In particular, we have proposed SIA, which links conditional answer entropy to the accumulation of answer-relevant information. SIA is not intended as a surprising claim; rather, it isolates the minimal structural condition under which entropy-based reasoning methods are theoretically justified, a condition that many empirical approaches in the literature implicitly rely on. Additionally, through a suite of experiments, we have verified that standard training pipelines induce model behavior consistent with SIA. We further found that correct reasoning traces exhibit characteristic entropy signatures that distinguish them from traces leading to incorrect answers with respect to the ground-truth distribution.

Lastly, some open questions remain. Entropy-based diagnostics may fail in regimes where reasoning-trace prefixes are only weakly informative about the true answer: characterizing the distributions that produce such behavior would clarify the limits of entropy as a proxy for reasoning. Also, it remains open whether targeted interventions that modify entropy dynamics can reliably change reasoning outcomes. Finally, an important direction is to generalize entropy-based diagnostics to other modalities and generative modeling paradigms.

\section*{Acknowledgements}
Mar Gonzàlez I Català acknowledges that this project was supported by G-Research.

\section*{Impact Statement}
This paper aims to advance the field of Machine Learning. While our work has potential societal implications, we do not identify any specific concerns that require particular emphasis at this stage.


\bibliographystyle{icml2026}
\bibliography{references}

\appendix
\onecolumn
\section{Experimental setup and evaluation protocol}
\label{app:exp}

\subsection{Evaluation protocol}
\label{subsec:eval_protocol}

\subsubsection{Tasks and datasets}
\label{subsec:tasks_datasets}

We focus on reasoning tasks with a discrete answer space \(\mathcal{A}\), which enables empirical estimation of conditional answer entropy. Each example consists of a question \(Q \in \mathcal{Q}\) and a ground-truth answer \(A \in \mathcal{A}\). We evaluate on the following datasets:

\begin{itemize}
\item \textbf{GSM8K} \cite{cobbe2021trainingverifierssolvemath}: grade-school mathematical word problems with numeric answers.
\item \textbf{ARC} \cite{clark2018thinksolvedquestionanswering}: multiple-choice science questions.
\item \textbf{SVAMP} \cite{patel2021nlpmodelsreallyable}: arithmetic word problems designed to test robustness to linguistic variation.
\end{itemize}

For all datasets, we use the official test splits and apply deterministic answer normalization and parsing to map model outputs to discrete answer labels (e.g., numeric normalization for GSM8K and SVAMP, letter-to-option mapping for ARC). Invalid or unparsable outputs are mapped to a special null answer category.

\subsubsection{Models}
\label{subsec:models}

We evaluate a diverse set of open-weight LLMs corresponding to different training regimes:

\begin{itemize}
\item \textbf{Gemma-2-2B} \cite{gemmateam2024gemma2improvingopen}: base and instruction-tuned variants.
\item \textbf{LLaMA-3.2-3B} \cite{grattafiori2024llama3herdmodels}: base and instruction-tuned variants.
\item \textbf{Qwen-2.5-3B} \cite{qwen2025qwen25technicalreport}: base and instruction-tuned variants.
\item \textbf{Qwen-2.5-Math-1.5B} \cite{qwen2025qwen25technicalreport}: SFT-trained specialized on math problems.
\item \textbf{DeepSeek-Chat-7B} \cite{deepseekai2024deepseekllmscalingopensource}: SFT-trained chat model.
\item \textbf{DeepSeek-R1-distilled-7B} \cite{Guo_2025}: reasoning-specialized RL model.
\item \textbf{Olmo-3-7B-Think} \cite{olmo2025olmo3}: SFT and RL-trained variants. 
\end{itemize}

Base models correspond to pretrained LLMs without supervised or reinforcement fine-tuning. Instruction-tuned (IT) models are supervised fine-tuned on instruction-following data. RL-trained models are optimized using reinforcement learning from human or synthetic feedback.

All models are evaluated using their publicly released checkpoints with default tokenizers and architectures.

\subsubsection{Generation procedure}
\label{subsec:generation_procedure}

For each question \(Q=q\), we sample \(M\) independent reasoning trajectories from the model under a fixed stochastic decoding configuration (temperature, nucleus sampling, and maximum generation length). Concretely, for each \(i \in \{1,\dots,M\}\) we draw
\[
C^{(i)}_{1:K^{(i)}} \sim p_\theta(\cdot \mid q),
\]
where \(K^{(i)}\) denotes the generated reasoning length (up to a fixed truncation limit). We treat each sampled trajectory \(C^{(i)}_{1:K^{(i)}}\) as one realization of the model's reasoning process for the given query.

Unless otherwise specified, decoding uses:
\begin{itemize}
    \item temperature $T=0.7$
    \item nucleus sampling with $p=0.9$
    \item a maximum generation length of 600 tokens
\end{itemize}

Each trajectory is treated as one realization of the model’s reasoning process for the given query. All rollouts used for entropy estimation use the same decoding configuration to ensure comparability.


\subsubsection{Monte-Carlo estimation of conditional answer entropy}
\label{subsubsec:mc_conditional_entropy}

Given a fixed query \(Q = q\) and a realized reasoning prefix \(C_{1:k} = c_{1:k}\), the model induces an implicit distribution over final answers
\[
p_\theta(A \mid q, c_{1:k}),
\]
We approximate \(H_{p_\theta}(A \mid q, c_{1:k})\) using Monte-Carlo sampling. For a fixed prefix \((q, c_{1:k})\), we draw \(N\) independent stochastic rollouts from the model:
\[
A^{(i)} \sim p_\theta(\cdot \mid q, c_{1:k}), \qquad i = 1, \dots, N.
\]
using the same decoding parameters as the base generation, followed by deterministic answer extraction.

These samples induce an empirical distribution
\[
\hat{p}_k(a \mid q, c_{1:k})
= \frac{1}{N} \sum_{i=1}^N \mathbf{1}\{A^{(i)} = a\}.
\]
and the plug-in estimator
\[
\widehat{H}_{p_\theta}(A \mid q, c_{1:k})
= -\mkern-46mu \sum_{a \in \mathcal{A} : \hat{p}_k(a \mid q, c_{1:k}) > 0}
\mkern-46mu 
\hat{p}_k(a \mid q, c_{1:k})
\log \hat{p}_k(a \mid q, c_{1:k}).
\]

This estimator is biased for finite \(N\) but consistent as \(N \to \infty\), and sufficient for comparing entropy trends across token positions and training regimes.

All rollouts are performed in evaluation mode, without gradient computation. Sampling parameters are held fixed across models and prefixes. In practice, we use $N=16$ continuations per prefix unless otherwise stated. For an ablation using $N=32$ continuations, see Appendix~\ref{subsec:further_ablations}.

\subsubsection{Checkpointed prefix evaluation}
\label{subsec:checkpointed_prefix_evaluation}

Estimating conditional answer entropy at every token is computationally expensive. We therefore evaluate at checkpoint positions
\[
\mathcal{K} = \{k_1, k_2, \dots, k_m\} \subseteq \{0,1,\dots,K\},
\]
spaced uniformly at stride $s$, and always including the final prefix length of the trajectory (\(k_m = K\)). Here \(k=0\) corresponds to the empty prefix.

For each \(k \in \mathcal{K}\), we compute \(\widehat{H}_{p_\theta}(A \mid Q, C_{\le k})\) independently. When needed for visualization, we linearly interpolate entropy values between checkpoints, but all reported quantitative results are computed on \(\mathcal{K}\).

\subsubsection{Statistical reporting}
\label{subsec:statistics}

Unless otherwise specified, all reported curves show the mean across questions, with shaded regions denoting 95\% bootstrap confidence intervals computed over questions. For metrics such as AUC or average information gain, we report both mean and standard error.
\onecolumn
\section{Further results}
\label{app:further_results}

\subsection{Signatures vanish or weaken in non-aligned models}
\label{subsec:signatures_weaken}

This appendix supports the claims made in Section~\ref{sec:signatures} that the observable signatures of Stepwise Informativeness (SIA) are specific to aligned models and either vanish or
significantly weaken in non-aligned ones.

\paragraph{Failure of early information accumulation.}
Figure~\ref{fig:nonaligned-earlyinfo} reports the normalized cumulative information gain
$G(s)$ for non-aligned models, split by correctness. Unlike aligned models (Figure~\ref{fig:auc_by_r}), correct traces do not exhibit systematically earlier or steeper accumulation of answer-relevant information. The two curves largely overlap, indicating the absence of early lock-in behavior.

\begin{figure}[hbtp!]
  \centering
  \includegraphics[width=0.55\linewidth]{figures/normalized_gain_curve_base.png}
  \caption{\textbf{Early information accumulation in non-aligned models.}
Normalized cumulative gain $G(s)$ vs.\ relative prefix length $s$, split by correctness in llama-3.2-3B (non-aligned model) in GSM8k dataset. Entropy is not a correctness signal in this regime.}
  \label{fig:nonaligned-earlyinfo}
\end{figure}

\paragraph{Failure of early separability.}
Figure~\ref{fig:nonaligned-separability} shows the AUC obtained when using conditional
answer entropy at prefix length $s$ to distinguish correct from incorrect traces.
For non-aligned models, separability remains weak across the entire generation and does
not rise sharply at small $s$, in contrast with the behavior observed in aligned models
(Figure~\ref{fig:auc_by_r}).

\begin{figure}[hbtp!]
  \centering
  \includegraphics[width=0.55\linewidth]{figures/separability_auc_by_r_base.png}
  \caption{\textbf{Separability in non-aligned models}.
  AUC for using conditional answer entropy to distinguish correct from incorrect traces vs.\ relative prefix length $s$ across non-aligned models in GSM8k dataset. Entropy is not an early diagnostic signal in this regime.}
  \label{fig:nonaligned-separability}
\end{figure}

Together, these results confirm that the empirical signatures described in Section~\ref{sec:signatures}
are not generic properties of autoregressive models, but arise specifically when training
induces stepwise informativeness.


\subsection{Further ablations}
\label{subsec:further_ablations}

\paragraph{Monte-Carlo approximation.}
Our entropy estimates rely on Monte-Carlo rollouts. To assess robustness to approximation quality, we reran a subset of experiments using a coarser estimator with stride $4$ and $32$ samples, on $100$ GSM8K instances across a subset of models. Table~\ref{tab:mc_ablation} reproduces a subset of Table~\ref{tab:sia_alignment_by_dataset} under this setting. Results remain qualitatively unchanged, indicating that SIA alignment is not an artifact of low-fidelity Monte-Carlo estimation.

\begin{table}[hbtp!]
\centering
\caption{Monte-Carlo ablation on GSM8K (stride $4$, MC=$32$, $100$ samples).}
\label{tab:mc_ablation}
\small
\begin{tabular}{lcc}
\toprule
\textbf{Model} & {Original mean} & {Ablated mean} \\
\midrule
Qwen2.5-3B      &  0.682 & 0.635 \\
Qwen2.5-3B-it   &  0.744 & 0.831 \\
DeepSeek-R1-Distilled    &  0.795 &  0.711 \\
Gemma-2-2B-it & 0.522 & 0.506 \\
\bottomrule
\end{tabular}
\end{table}

\onecolumn
\section{Proofs}
\label{app:proofs}

\begin{proof}[Proof of Lemma~\ref{cl:exp_delta_is_I}]
The expectation of $\Delta_k$ expands as
\begin{equation*}
\mathbb{E}[\Delta_k] = \sum_{q,a,c_{1:k}} p(q,a,c_{1:k})\, \log\frac{p(a \mid q,c_{\le k})}{p(a \mid q,c_{<k})}.
\end{equation*}

To make the relationship with conditional mutual information explicit, we separate the prefix $c_{<k}$ from the current token $c_k$:
\begin{equation*}
\mathbb{E}[\Delta_k]
=
\sum_{q,c_{<k},c_k,a}
p(q,c_{<k},c_k,a)\,
\log\frac{p(a \mid q,c_{<k},c_k)}{p(a \mid q,c_{<k})}.
\end{equation*}

We can rewrite the expectation in terms of the conditional distribution of $(A,C_k)$ given $(Q,C_{<k})$:
\begin{equation*}
\mathbb{E}[\Delta_k] = \sum_{q,c_{<k}} p(q,c_{<k}) \sum_{c_k,a} p(a,c_k \mid q,c_{<k}) \log\frac{p(a \mid q,c_{<k},c_k)}{p(a \mid q,c_{<k})}.
\end{equation*}

Using the factorization
\begin{equation*}
p(a,c_k \mid q,c_{<k})
= p(a \mid q,c_k,c_{<k})\, p(c_k \mid q,c_{<k}),
\end{equation*}
we recognize that by rewriting the logarithm inside the sum we obtain exactly the definition of the conditional mutual information:
\begin{equation*}
\mathbb{E}[\Delta_k] = \sum_{q,c_{<k}} p(q,c_{<k}) \sum_{c_k,a} p(a,c_k \mid q,c_{<k}) \log\frac{p(a,c_k \mid q,c_{<k})}{p(a \mid q,c_{<k})\,p(c_k \mid q,c_{<k})} = I(A;C_k \mid Q,C_{<k}).
\end{equation*}

Also, we can express the mutual information in terms of entropy. 
\begin{equation*}
I(A;C_k \mid Q,C_{<k}) = \sum_{q,c_{<k},c_k,a}
p(q,c_{<k},c_k,a)\,\log\frac{p(a,c_k \mid q,c_{<k})}{p(a \mid q,c_{<k})\,p(c_k \mid q,c_{<k})}
\end{equation*}

\begin{equation*}
= \sum_{q,c_{<k},c_k,a}
p(q,c_{<k},c_k,a)\,\log\frac{p(a\mid q,c_{<k},c_k)p(c_k \mid q,c_{<k})}{p(a \mid q,c_{<k})\,p(c_k \mid q,c_{<k})}
\end{equation*}
\begin{equation*}
= \sum_{q,c_{<k},c_k,a}
p(q,c_{<k},c_k,a)\,\big(\log p(a\mid q,c_{<k},c_k) - \log p(a \mid q,c_{<k})\big)
\end{equation*}

\begin{equation*}
= - \sum_{q,c_{<k},c_k,a}
p(q,c_{<k},c_k,a)\log p(a \mid q,c_{<k}) +\sum_{q,c_{<k},c_k,a}
p(q,c_{<k},c_k,a)\,\log p(a\mid q,c_{\le k}) 
\end{equation*}

Next, consider the first term

\begin{equation*}
-\sum_{q,c_{<k},c_k,a} p(q,c_{<k},c_k,a) \log p(a \mid q,c_{<k}).
\end{equation*}

Notice that the probability $p(a \mid q,c_{<k})$ does not depend on $c_k$. Therefore we can rewrite the sum as

\begin{equation*}
-\sum_{q,c_{<k},a} \log p(a \mid q,c_{<k}) \left( \sum_{c_k} p(q,c_{<k},c_k,a) \right).
\end{equation*}

The inner sum is simply the marginal probability obtained by summing over $c_k$:

\begin{equation*}
\sum_{c_k} p(q,c_{<k},c_k,a) = p(q,c_{<k},a).
\end{equation*}

Substituting this back in gives

\begin{equation*}
H(A \mid Q, C_{<k}):=-\sum_{q,c_{<k},a} p(q,c_{<k},a) \log p(a \mid q,c_{<k}).
\end{equation*}

Hence, 
\begin{equation*}
I(A;C_k \mid Q,C_{<k}) = H(A \mid Q,C_{<k}) - H(A \mid Q,C_{\le k}),
\end{equation*}

and we arrive at the compact form
\begin{equation*}
\mathbb{E}[\Delta_k] = I(A;C_k \mid Q,C_{<k}) = H(A\mid Q,C_{<k}) - H(A\mid Q,C_{\le k}).
\end{equation*}
\end{proof}

\begin{proof}[Proof of Proposition~\ref{prop:entropy-information}]
For each $t \ge 1$, the conditional mutual information admits the standard entropy decomposition
\[
I(A; C_t \mid Q, C_{<t})
=
H(A \mid Q, C_{<t}) - H(A \mid Q, C_{\le t}).
\]
Summing over $t = 1,\dots,k$ yields a telescoping series:
\begin{align*}
\sum_{t=1}^k I(A; C_t \mid Q, C_{<t})
&=
\sum_{t=1}^k \bigl[ H(A \mid Q, C_{<t}) - H(A \mid Q, C_{\le t}) \bigr] \\
&=
H(A \mid Q) - H(A \mid Q, C_{1:k}),
\end{align*}
which establishes the first identity.

The second identity follows from the chain rule for conditional mutual information, which states that
\[
I(A; C_{1:k} \mid Q)
=
\sum_{t=1}^k I(A; C_t \mid Q, C_{<t}).
\]
Combining the two expressions completes the proof.
\end{proof}

\begin{proof}[Proof of Theorem~\ref{th:entropy_accuracy}]
Consider the pair of random variables
\[
A \in \mathcal{A} \qquad Y := (Q, C_{1:k})
\]
with $|\mathcal A|\ge 2$. Let $\hat A(Y)$ be the Bayes-optimal (MAP) classifier under $p(A\mid Y)$ and denote
\[
P_e(Y) := \Pr(\hat A(Y)\neq A).
\]
Fano's inequality \cite{Cover2006} states that 
\[
H(A\mid Y) \le \log 2 + P_e(Y)\log(|\mathcal A|-1),
\]
which rearranges to
\[
P_e(Y) \ge \frac{H(A\mid Y)-\log 2}{\log(|\mathcal A|-1)}.
\]
Substituting $Y=(Q,C_{1:k})$ yields
\[
P_e^{(k)} \ge \frac{H(A\mid Q,C_{1:k})-\log 2}{\log(|\mathcal A|-1)}.
\]
\end{proof}



\begin{proof}[Proof of Lemma~\ref{lemma:log-likelihood_identity}]
By definition of expectation under $r$, we have
\[
\mathcal{L}(\theta)
= \mathbb{E}_{X\sim r}[-\log p_\theta(X)]
= -\sum_{x} r(x)\log p_\theta(x).
\]
We now add and subtract $\sum_{x} r(x)\log r(x)$, which equals zero:
\[
\mathcal{L}(\theta)
= -\sum_{x} r(x)\log p_\theta(x)
  + \sum_{x} r(x)\log r(x)
  - \sum_{x} r(x)\log r(x).
\]
Rearranging the terms gives
\[
\mathcal{L}(\theta)
= -\sum_{x} r(x)\log r(x)
  + \sum_{x} r(x)\log\frac{r(x)}{p_\theta(x)}.
\]
The first term is the Shannon entropy
\[
H(r) = -\sum_{x} r(x)\log r(x),
\]
which depends only on the data-generating distribution $r$ and not on $\theta$. It represents the irreducible uncertainty of the data source. In natural language, this idea goes back to Shannon’s analysis of the entropy of printed English \citep{6773263}, and in modern language modeling manifests as a non-zero lower bound on achievable cross-entropy or perplexity, as observed in empirical scaling laws \citep{kaplan2020scalinglawsneurallanguage}. 

The second term is the forward Kullback--Leibler divergence
\[
\mathrm{KL}(r \,\Vert\, p_\theta)
= \sum_{x} r(x)\log \frac{r(x)}{p_\theta(x)}.
\]
Thus we obtain the exact decomposition
\[
\mathcal{L}(\theta)
= H(r) + \mathrm{KL}(r \,\Vert\, p_\theta).
\]

Since $H(r)$ is constant with respect to~$\theta$, minimizing $\mathcal{L}(\theta)$ is equivalent to minimizing $\mathrm{KL}(r \,\Vert\, p_\theta)$.   Therefore, any sequence of parameters $\theta$ that decreases the negative log-likelihood necessarily drives the model distribution $p_\theta$ toward the data distribution $r$ in Kullback--Leibler divergence.  This establishes that $p_\theta \approx r$ whenever $\mathcal{L}(\theta)$ is near its minimum.
\end{proof}

\begin{proof}[Proof of Lemma~\ref{lemma:KLDecompositionoftheJointConditional}]
Using the chain rule of probability, both distributions factorize as
\[
r(C_{1:K},A\mid Q)
= r(C_{1:K}\mid Q)\, r(A\mid Q,C_{1:K}),
\qquad
p_\theta(C_{1:K},A\mid Q)
= p_\theta(C_{1:K}\mid Q)\, p_\theta(A\mid Q,C_{1:K}).
\]
Hence the KL divergence expands to
\[
\mathrm{KL}
= 
\mathbb{E}_{r(C_{1:K},A\mid Q)}
\left[
\log \frac{
r(C_{1:K}\mid Q)\, r(A\mid Q,C_{1:K})
}{
p_\theta(C_{1:K}\mid Q)\, p_\theta(A\mid Q,C_{1:K})
}
\right].
\]
Splitting the logarithm into two terms yields
\[
\mathrm{KL}
=
\mathbb{E}_{r(C_{1:K},A\mid Q)}
\!\left[
\log\frac{r(C_{1:K}\mid Q)}{p_\theta(C_{1:K}\mid Q)}
\right]
+
\mathbb{E}_{r(C_{1:K},A\mid Q)}
\!\left[
\log\frac{r(A\mid Q,C_{1:K})}{p_\theta(A\mid Q,C_{1:K})}
\right].
\]
In the first term the integrand depends only on $C_{1:K}$, so the outer expectation reduces to
$\mathbb{E}_{r(C_{1:K}\mid Q)}$, giving
\[
\mathbb{E}_{r(C_{1:K}\mid Q)}
\left[
\log\frac{r(C_{1:K}\mid Q)}{p_\theta(C_{1:K}\mid Q)}
\right]
= 
\mathrm{KL}\!\left(r(C_{1:K}\mid Q)\,\Vert\,p_\theta(C_{1:K}\mid Q)\right).
\]
For the second term, conditioning on $C_{1:K}$ gives
\[
\mathbb{E}_{r(C_{1:K}\mid Q)}
\!\Big[
\mathbb{E}_{r(A\mid Q,C_{1:K})}
\left[
\log\frac{r(A\mid Q,C_{1:K})}{p_\theta(A\mid Q,C_{1:K})}
\right]
\Big]
=
\mathbb{E}_{r(C_{1:K}\mid Q)}
\!\left[
\mathrm{KL}\!\left(r(A\mid Q,C_{1:K}) \,\Vert\, p_\theta(A\mid Q,C_{1:K})\right)
\right].
\]
Combining both parts yields the claimed identity.
\end{proof}

\begin{proof}[Proof of Lemma~\ref{lemma:MLE_implies_alignment}]
By the decomposition in Lemma~\ref{lemma:joint-kl-decomposition},
the joint KL is a sum of two nonnegative terms:
\[
\mathrm{KL}(r\Vert p_\theta)
=
\underbrace{
\mathrm{KL}\!\left(r(C_{1:K}\mid Q)\Vert p_\theta(C_{1:K}\mid Q)\right)
}_{\text{marginal term}}
+
\underbrace{
\mathbb{E}_{r(C_{1:K}\mid Q)} 
\left[
\mathrm{KL}\!\left(r(A\mid Q,C_{1:K})\Vert p_\theta(A\mid Q,C_{1:K})\right)
\right]
}_{\text{conditional term}}.
\]
Thus, if the sum is bounded by $\delta$, then each individual term must also be bounded by
$\delta$:
\[
\mathrm{KL}\!\left(
r(C_{1:K}\mid Q)
\Vert
p_\theta(C_{1:K}\mid Q)
\right)
\le \delta,
\]
and
\[
\mathbb{E}_{r(C_{1:K}\mid Q)}
\!\left[
\mathrm{KL}\!\left(
r(A\mid Q,C_{1:K})
\Vert
p_\theta(A\mid Q,C_{1:K})
\right)
\right]
\le \delta.
\]

This establishes both claims.
\end{proof}

\begin{proof}[Proof of Lemma~\ref{lemma:entropy-continuity}]
Let $\|\cdot\|_{\mathrm{TV}}$ denote total variation distance,
\[
\|P - Q\|_{\mathrm{TV}} := \frac12 \sum_{x \in \mathcal{X}} |P(x) - Q(x)|.
\]
By Pinsker's inequality,
\[
\|P - Q\|_{\mathrm{TV}} \le \sqrt{\frac12 \mathrm{KL}(P \Vert Q)}
\le \sqrt{\frac{\delta}{2}}.
\]
Let $\varepsilon := \|P-Q\|_{\mathrm{TV}}$. The Fannes--Audenaert inequality \cite{Audenaert_2007} (continuity of entropy on a finite alphabet) states that
for $\varepsilon \le 1 - \tfrac{1}{|\mathcal{X}|}$,
\[
\bigl|H(P) - H(Q)\bigr|
\le \varepsilon \log(|\mathcal{X}| - 1) + h_2(\varepsilon),
\]
where $h_2(\varepsilon) := -\varepsilon \log \varepsilon
- (1-\varepsilon)\log(1-\varepsilon)$ is the binary entropy function.

Combining these two inequalities, for all $\delta > 0$ such that
$\sqrt{\delta/2} \le 1 - 1/|\mathcal{X}|$ we obtain
\[
\bigl|H(P) - H(Q)\bigr|
\le f_\mathcal{X}(\delta),
\]
where one admissible choice is
\[
f_\mathcal{X}(\delta)
:= \sqrt{\frac{\delta}{2}} \log(|\mathcal{X}| - 1)
+ h_2\!\Bigl(\sqrt{\tfrac{\delta}{2}}\Bigr).
\]
The function $f_\mathcal{X}$ is continuous and satisfies
$f_\mathcal{X}(\delta) \to 0$ as $\delta \to 0$, because both
terms on the right-hand side vanish in this limit.

Finally, the $\varepsilon$--$\delta$ formulation follows by continuity:
for any fixed $\varepsilon > 0$, choose $\delta > 0$ such that
$f_\mathcal{X}(\delta) \le \varepsilon$.
\end{proof}

\begin{proof}[Proof of Lemma~\ref{lemma:cond-entropy-continuity}]
Recall that conditional entropy can be expressed in terms of joint
entropies:
\[
H_P(Y \mid X) = H_P(X,Y) - H_P(X), \qquad
H_Q(Y \mid X) = H_Q(X,Y) - H_Q(X).
\]
Therefore
\begin{align*}
\bigl| H_P(Y \mid X) - H_Q(Y \mid X) \bigr|
&= \bigl| \bigl(H_P(X,Y) - H_Q(X,Y)\bigr)
- \bigl(H_P(X) - H_Q(X)\bigr) \bigr| \\
&\le \bigl|H_P(X,Y) - H_Q(X,Y)\bigr|
+ \bigl|H_P(X) - H_Q(X)\bigr|.
\end{align*}

Let $P_{XY}$ and $Q_{XY}$ denote the joint distributions on
$\mathcal{X}\times\mathcal{Y}$, and $P_X$, $Q_X$ the corresponding
marginals on $\mathcal{X}$. Since marginalization cannot increase
KL divergence, we have
\[
\mathrm{KL}(P_X \Vert Q_X)
\le \mathrm{KL}(P_{XY} \Vert Q_{XY})
= \mathrm{KL}(P\Vert Q) \le \delta.
\]
Applying Lemma~\ref{lemma:entropy-continuity} first to
$(P_{XY},Q_{XY})$ on the alphabet $\mathcal{X}\times\mathcal{Y}$ and
then to $(P_X,Q_X)$ on the alphabet $\mathcal{X}$ yields
\[
\bigl|H_P(X,Y) - H_Q(X,Y)\bigr|
\le f_{\mathcal{X}\times\mathcal{Y}}(\delta), \qquad
\bigl|H_P(X) - H_Q(X)\bigr|
\le f_{\mathcal{X}}(\delta).
\]
Combining the bounds, we obtain
\[
\bigl| H_P(Y \mid X) - H_Q(Y \mid X) \bigr|
\le f_{\mathcal{X}\times\mathcal{Y}}(\delta)
+ f_{\mathcal{X}}(\delta)
=: g_{\mathcal{X},\mathcal{Y}}(\delta),
\]
where $g_{\mathcal{X},\mathcal{Y}}(\delta)\to 0$ as $\delta\to 0$
because each $f$ has this property.
The $\varepsilon$--$\delta$ formulation follows as before.
\end{proof}



\begin{proof}[Proof of Lemma~\ref{lemma:mi-continuity}]
We work with finite alphabets, so all random variables take values in finite sets. Recall the standard entropy identity
\[
I(A; C_{\le k} \mid Q)
=
H(A \mid Q)
-
H(A \mid Q, C_{\le k}).
\]
For the distribution $r$ we have
\[
I_r(A; C_{\le k} \mid Q)
=
H_r(A \mid Q)
-
H_r(A \mid Q, C_{\le k}),
\]
and similarly for $p_\theta$:
\[
I_{p_\theta}(A; C_{\le k} \mid Q)
=
H_{p_\theta}(A \mid Q)
-
H_{p_\theta}(A \mid Q, C_{\le k}).
\]
Subtracting the two expressions gives
\begin{align*}
&\bigl|
I_r(A; C_{\le k} \mid Q)
-
I_{p_\theta}(A; C_{\le k} \mid Q)
\bigr| \\
&\qquad\le
\bigl|
H_r(A \mid Q)
-
H_{p_\theta}(A \mid Q)
\bigr|
+
\bigl|
H_r(A \mid Q, C_{\le k})
-
H_{p_\theta}(A \mid Q, C_{\le k})
\bigr|.
\end{align*}

We now apply Lemma~\ref{lemma:cond-entropy-continuity} twice, once
with $(X,Y) = (Q,A)$ and once with $(X,Y) = (Q,C_{\le k},A)$. Since marginalization cannot increase KL divergence, we have
\[
\mathrm{KL}\bigl(r(Q,A)
\Vert
p_\theta(Q,A)
\bigr)
\le \mathrm{KL}(r\Vert p_\theta)
\le \delta,
\]
and similarly for the joint $(Q,C_{\le k},A)$. Thus, Lemma
\ref{lemma:cond-entropy-continuity} yields functions
$g^{(<k)}(\delta)$ and $g^{(\le k)}(\delta)$, each vanishing as
$\delta\to 0$, such that
\[
\bigl|
H_r(A \mid Q)
-
H_{p_\theta}(A \mid Q)
\bigr|
\le g^{(<k)}(\delta),
\]
and
\[
\bigl|
H_r(A \mid Q, C_{\le k})
-
H_{p_\theta}(A \mid Q, C_{\le k})
\bigr|
\le g^{(\le k)}(\delta).
\]
Combining these inequalities gives
\[
\bigl|
I_r(A; C_{\le k} \mid Q)
-
I_{p_\theta}(A; C_{\le k} \mid Q)
\bigr|
\le
g^{(<k)}(\delta) + g^{(\le k)}(\delta)
=: G_k(\delta),
\]
where $G_k(\delta)\to 0$ as $\delta\to 0$. The
$\varepsilon$--$\delta$ formulation is immediate.
\end{proof}

\begin{proof}[Proof of Theorem~\ref{thm:transfer-sia}]
Given a step $k \geq 1$, we have
\[
I_r(A; C_{\leq k} \mid Q) > 0.
\]
Then there exists $\varepsilon_k > 0$ such that
\[
I_r(A; C_{\leq k} \mid Q) \ge \varepsilon_k > 0.
\]
Lemma~\ref{lemma:mi-continuity} (continuity of conditional mutual information) states that if 
\[
\mathrm{KL}(r \Vert p_\theta) \le \delta,
\]
then there exists a function $G_k(\delta)$ with $G_k(\delta) \to 0$ as $\delta \to 0$ such that
\[
\bigl| I_r(A; C_{\le k} \mid Q) - I_{p_\theta}(A; C_{\le k} \mid Q) \bigr|
\;\le\; G_k(\delta).
\]
Equivalently,
\[
I_{p_\theta}(A; C_{\le k} \mid Q)
\ge I_r(A; C_{\le k} \mid Q) - G_k(\delta).
\]
Combining this inequality with SIA we have
\[
I_{p_\theta}(A; C_{\le k} \mid Q)
\ge \varepsilon_k - G_k(\delta).
\]
If $\delta$ is chosen such that $G_k(\delta) < \varepsilon_k/2$, then
\[
I_{p_\theta}(A; C_{\le k} \mid Q)
\ge \varepsilon_k - G_k(\delta)
\ge \frac{\varepsilon_k}{2}.
\]
This proves the approximate SIA inequality for the model at step $k$.
\end{proof}

\end{document}
